\chapter{한계}
\label{ch:limitations}

본 장의 항목들은 조용히 고칠 결함이 아니라 방법 또는 현재 연구 상태의 성질이다.

\section{방법의 한계}

\paragraph{조각마다 상수 가정.}
기본 Chan--Vese는 내부와 외부의 밝기가 각각 일정하다는 단순한 가정을 사용한다.
Papillary muscle, blood pool, myocardium이 섞인 영역에서는 국소 영역 항이나 multiphase
model이 필요할 수 있다. 본 연구의 팬텀은 이 위반을 의도적으로 포함하지만, 그것을
해결하지는 않는다.

\paragraph{접선 운동의 미결정성.}
법선 정사영은 식 \eqref{eq:least-norm}의 최소 노름 해를 택하므로 접선 방향 성분을
결정하지 않는다. 따라서 결과는 기하학적 표면 운동이며 실제 조직 대응이나 strain이
아니다. Vertex correspondence가 중요한 기능적 계측에는 mesh 운동 추적이 더 적합하다.

\paragraph{위상 변화와 근접 분기.}
정사영은 인접 프레임 사이의 이동이 작고 동일한 표면 분기가 유지될 때 안정적이다.
큰 이동, 높은 곡률, 가까운 두 경계, 위상 변화에서는 잘못된 표면으로 이동할 수 있다.
Trust region과 되돌림이 이를 줄이지만 제거하지는 못한다.

\paragraph{표면 렌더러의 범위.}
2DGS는 명시적 표면 렌더러이다. 심장 내부 전체의 정확한 volumetric MRI나 임의 단면을
복원하려면 별도 volumetric model이 필요하며, 이는 본 연구의 주장이 아니다.

\paragraph{잔차 저장량.}
진폭 잔차를 surfel마다 모든 프레임에 저장하면 surfel 수가 매우 클 때 저장량이
증가한다. 식 \eqref{eq:lowrank}의 저계수 압축이 이를 완화하지만, 명제
\ref{prop:lowrank}은 오차가 특잇값 꼬리와 같음을 말할 뿐 꼬리가 작음을 보장하지 않는다.

\paragraph{입력 모달리티의 차이.}
원 2DGS는 calibrated multi-view RGB 영상을 가정하지만 cine CMR은 물리 좌표가 알려진
여러 MRI slice이다. 본 연구는 2DGS를 새로운 장면복원기로 쓰지 않고 알려진 geometry로
투영 가능한 표면 렌더러로 사용한다. 감독도 알려진 slice plane의 네 신호로 제한된다.

\paragraph{원근 정합의 계산 비용.}
화소별 광선--surfel 교차는 단순 삼각형 래스터화보다 primitive당 계산이 더 클 수 있다.
따라서 실시간성과 mesh 대비 품질 이득은 \emph{별개의} 가설이며 둘 다 실제 하드웨어에서
측정해야 한다.

\paragraph{중간 시점 보간.}
명제 \ref{prop:interp-eikonal}, \ref{prop:interp-projection}에 의해 보간된 level-set은
참 중간 표면의 부호 거리 함수가 아니다. 표시 전용이며 임상적 복원이 아니다.

\section{범위의 한계}

\paragraph{단일 구조.}
좌심실 cavity만 구현되었다. 전체 좌/우심실 및 심근은 multiphase level-set을 요구하는
확장 범위이며, 실제 데이터 로더는 다른 구조를 요청하면 예외를 발생시킨다.

\paragraph{Digital twin이 아니다.}
조직 물성, 전기전도, 혈류, 역학 모델을 포함하지 않으며 치료 반응을 예측하지 않는다.
결과물은 시간가변 해부학 시각화 모델이다.

\paragraph{정답의 부재.}
ACDC와 M\&Ms-2는 ED/ES에만 label을 제공한다. 중간 프레임에는 분할 정답이 없으므로
그 프레임들은 외형과 시간적 안정성 지표로만 평가할 수 있다.

\section{현재 연구 상태의 한계}

\caveat{%
\textbf{구현이 실행되지 않았다.} 본 논문에서 가장 중대한 한계이다. PyTorch를 설치할 수
없는 환경에서 작성되었으므로 RQ1--RQ6의 어떤 주장도 확인되지도 반박되지도 않았다.
문법, import 해소, 이산 논리 커널 41개 항목은 검증되었으나(제
\ref{sec:kernel-verification}절) 그것은 파이프라인이 동작함을 뜻하지 않는다. 첫 실행에서
결함이 드러날 것으로 예상하는 것이 합리적이며, 그 목적으로 21스테이지 스모크 진단을
함께 제공하였다.}

\paragraph{통계가 구현되지 않았다.}
제 \ref{sec:statistics}절이 요구하는 환자 수준 paired 통계(bootstrap 신뢰구간,
Wilcoxon signed-rank)는 구현되지 않았다. 지표는 환자별로 산출되지만 집계는 실제
데이터와 함께 추가되어야 한다. 단일 팬텀 결과는 환자에 대한 증거로 읽을 수 없다.

\paragraph{비열등성 마진이 정해지지 않았다.}
결과를 본 뒤에 정하면 검정이 무의미해지므로 의도적으로 비워 두었다. pilot이 선행되어야
한다.

\paragraph{Dyna3DGR이 재현되지 않았다.}
가장 가까운 선행연구이지만 재현하지 않았고 어떤 수치도 귀속시키지 않았다. 구조적
비교만 제시하였다(표 \ref{tab:dyna3dgr}).

\paragraph{실제 데이터가 실행되지 않았다.}
ACDC/M\&Ms-2 로더는 구현되었으나 데이터에 접근할 수 없어 실행되지 않았다. 따라서
실제 해부학, scanner 변이, 질환에 대한 강건성은 전혀 평가되지 않았다.

\paragraph{PDF가 컴파일되지 않았다.}
본 원고를 작성한 환경에는 \LaTeX{} 배포판이 없었다. 따라서 이 문서 자체도 컴파일되어
확인된 적이 없으며, 조판 오류가 남아 있을 수 있다.

\section{다음에 해야 할 일}

우선순위 순서이다.

\begin{enumerate}[leftmargin=2em]
  \item PyTorch가 있는 환경에서 \texttt{python scripts/verify\_kernels.py}로 41/41을
  재확인한 뒤 \texttt{cvdyn2dgs smoke}로 파이프라인을 올린다.
  \item \texttt{cvdyn2dgs theory}로 제 \ref{ch:error}장의 수렴률을 측정한다. 이것이
  실패하면 이후 품질 수치는 무의미하다.
  \item 팬텀에서 RQ1--RQ6을 실행하고 \texttt{scripts/make\_tables.py}로 제
  \ref{ch:results}장을 채운다.
  \item ACDC 5명으로 pilot을 수행하고 비열등성 마진을 고정한다.
  \item ACDC 20--30명, M\&Ms-2 10--20명으로 본 실험을 수행하고 환자 수준 paired 통계를
  추가한다.
  \item 제 \ref{ch:discussion}장의 사전 해석 규칙에 따라 결론을 작성한다.
\end{enumerate}
